\documentclass[a4paper,10pt]{article}
\newcommand\subdir{/home/koen/LaTeX-setup/}
\input{\subdir/LaTeX-files/imports.tex}
\input{\subdir/LaTeX-files/master-internship.tex}
\usepackage{\subdir/LaTeX-files/aastex}
\usepackage{natbib}
\bibliographystyle{\subdir/LaTeX-files/astroadstitle.bst}

%Set the title, Author and date
\title{Notes Week 31}
\author{Koen Essers}
\tutorial{Master Internship}
\date{\today}

%Set the header style
\header{\tutorialstring}

%Begin the document
\begin{document}
\setuplayout
\section{Fully Ionized Mean Molecular Weight}

\subsection{Main Sequence}

Since we are using the fully ionized mean molecular weights for the accreted material, it is only fair to compare this with the
fully ionized mean molecular weights of MESA abundances.

\begin{figure}[ht]
	\centering
	\incplot{w31-computed-mu}
\end{figure}

The gray lines are the \(\mu \) values as reported by \emph{MESA}.
The colored lines are the
computed values from the abundances and the horizontal line is \(\mu \) from \citet{asplund2009}.
I \emph{really} do not understand why the computed mean molecular weights from the \emph{MESA}
models are lower than reported \(\mu \) from \emph{MESA} and the computed \citet{asplund2009} \(\mu \)
values.

I feel like the \(\mu \) values reported from \emph{MESA} are more representative because they
agree with \citet{asplund2009}(?
)

\subsection{Giant branch and Thermally Pulsing Asymptotic Giant Branch}

\begin{figure}[ht]
	\centering
	\incplot{w31-GB-mu}
\end{figure}

\begin{figure}[ht]
	\marginfignote{0}{The mean molecular weights reported by the \emph{MESA} models on the main sequence agree
		\emph{very} well with the mean molecular weights of the \citet{asplund2009}.
	}
	\marginfignote{9}{
		On the TPAGB, it is exactly opposite.
		Here, the computed values
		from the abundances agrees much better with the \(\mu \) computed from the \citet{karakas2016} abundances.
	}
	\centering
	\incplot{w31-re}
\end{figure}

\newsection{Failing models II}

I wanted to look into the models that fail when undergoing RLOF early on in the evolution.

Below, I start by plotting the ratio between the star radius and the outer Roche lobe radius as
a function of the envelope mass.

\begin{figure}[ht]
	\marginfignote{0}{Clearly the models that fail all \emph{would} obtain very large outer lobe
		overfill ratios.
		However, this is clearly not \emph{the} reason that the models crash.
	}
	\centering
	\incplot{w31-ROL}
\end{figure}

I read \citet{temmink2022} again carefully and there is one thing that I think could be important.
If I understand the assumptions in the paper correctly, the idea is the following:

If the mass loss rate is lower than a critical mass loss rate \(\dot{M}_\text{crit}\), the local
thermal timescale of outer layers is able to ''keep up`` with the envelope stripping and it
can maintain its local thermal equilibrium.
This thermal equilibrium is also what allows the star
to maintain its superadiabatic surface layers.

\citet{temmink2022} then writes: ``In this scenario, the new sub-surface layers have a lower entropy than before the mass loss.
Thus, similar to what happens in a radiative star, the sub-surface layers occupy a smaller volume after decompression.
''

If I understand this correctly, any \(\delta S < 0\) for a layer would mean the radius of that layer decreases?
Is that right?

I noticed some interesting things in the entropy profiles of the stars.
Below, I plot the minimum entropy that the stars have in their envelope.
This is generally
at the end of the superadiabatic sub-surface layers.
After this minimum, the entropy slightly
increases again towards the surface.

\begin{figure}[ht]
	\marginfignote{0}{Unfortunately, I have a limited amount of models, so the lines are not
		very smooth, \emph{but}, it should be clear that the models that crash actually increase
		their minimum entropy when they are stripped.
		At some point, when the initial Roche lobe radius
		becomes large enough, this is no longer the case and the minimum entropy stays pretty constant.
	}
	\centering
	\incplot{w31-min-entropy}
\end{figure}

I decided to simulate a grid of models with initial Roche lobe radii between the last model that
failed at the first model that succeeded.
I saved a bunch of custom additional history data to
hopefully help diagnose this.

\begin{figure}[ht]
	\marginfignote{0}{This first plot simply shows the Roche lobe radius and the radius of the stars
		as a function of time.
		We can see that all of the failing models interact during 1 thermal pulse,
		while all of the succeeding models interact 1 thermal pulse later.
	}
	\marginfignote{20}{The bottom panel again shows the minimum entropy as a function of envelope mass.
		This time the resolution is much better because I compute this for every history step.
		We can now even more clearly see that the failing models tend to get much larger
		minimum entropies when losing their envelope.
	}
	\centering
	\incplot{w31-entropy-2}
\end{figure}

I also tried to create a similar plot for the entropy in the convective part of the envelope,
but my code for selecting this was not great.
Instead, I show the entropy profiles for a star that completed RLOF:
\begin{figure}[ht]
	\marginfignote{0}{Clearly, not only the minimum entropy changes, but the average entropy
		in the convective envelope changes too.
		I am not sure how this would change the
		mass transfer stability criterion.
		I think larger entropy should mean a more extended star, but
		I am not sure.
	}
	\centering
	\incplot{w31-full-S-profile}
\end{figure}

I then also looked at the width of the superadiabatic region.
\begin{figure}[ht]
	\marginfignote{0}{I am not super sure \emph{why} this
		would be important, but there is again a clear distinction between the models that fail and those that do not.
		In general, the larger initial Roche lobe radii have larger superadiabatic regions when losing mass.
		This is probably due to their later evolutionary state when interacting.
	}
	\centering
	\incplot{w31-dm-sad}
\end{figure}

The figure below shows the more detailed values

\begin{figure}[ht]
	\marginfignote{0}{The large jumps towards lower envelope masses are artificial and
		originate only from the wavey nature of the superadiabicity profile.
		I show this
		in the plot below.
	}
	\marginfignote{7}{We do see though that the superadiabatic region is smallest in the
		models that crash.
		From the plot above, it also looks like the relative
		width also increases quite a lot with the age of the TPAGB star, or the initial Roche
		lobe radius, which is very similar in this case.
	}
	\centering
	\incplot{w31-sad-m-2}
\end{figure}

\begin{figure}[ht]
	\marginfignote{0}{The horizontal line shows the criterion that \citet{temmink2022} uses for superadiabicity.
		The waves can cause models to show large jumps in the plot above.
		These are not "real".
	}
	\marginfignote{10}{This is a \(R_\text{RL,i =236.875}\; R_\odot \), \(M_\text{TAPGB,i} = 2.2\; M_\odot \), \(q=0.6\) and \(\varepsilon =0.25\) model.}
	\centering
	\incplot{w31-superadiabatic-1}
\end{figure}

I want to return to this plot I made a while ago:

\begin{figure}[ht]
	\marginfignote{0}{This is a grid with a \(M=2.0\; M_\odot \) TPAGB-star, \(\gamma =1\)
		and \(\delta =0.5\).
		The pink squares are models that fail due to convergence problems.
		The green models end because the models spiral in.
	}
	\centering
	\incplot{w15-grid-1}
\end{figure}

I am curious if the models that fail, all happen in during the same thermal pulses.
If however, it is possible to survive RLOF during a thermal pulse with a certain \(q\), but
not with another \(q\).

\begin{figure}[ht]
	\marginfignote{0}{Clearly there are models that interact at roughly the same time (same
		color), but some of them crash, (dashed) while others do not.}
	\centering
	\incplot{w31-grid-1}
\end{figure}
\nextpage

\begin{figure}[ht]
	\marginfignote{0}{Here, I explicitly show this.
		The models interact during the same pulse,
		however, only the green lines, with the largest \(q\), survive.
	}
	\centering
	\incplot{w31-see-same-pulse}
\end{figure}

I guess the mass loss rate is another important consideration.

\begin{figure}[ht]
	\marginfignote{0}{I show that here.
		The models with the largest \(q\) have the
		smallest \(\dot{M}\) and thus strip the star slower.
		The differences here might seem minor, but this is a log scale so differences are
		actually quite large.
	}
	\centering
	\incplot{w31-see-same-pulse-2}
\end{figure}

\begin{figure}[ht]
	\centering
	\incplot{w31-see-same-pulse-3}
\end{figure}

\section{Grids}

Using a very simple assumption, it is possible to quite accurately estimate the final accretor
abundances of the detailed \emph{MESA} simulations using only the simple direct binary evolution
in python with \lstinline[style=output, breaklines=true]|evolve.py|.
The idea is this:

Follow the simple binary evolution as predicted by \lstinline[style=output, breaklines=true]|evolve.py| exactly up to the point of RLOF.
Once RLOF happens, assume that \emph{all} of the remaining envelope mass is expelled instantly and accreted with a certain efficiency.
To match \emph{MESA} models, we should let \(\varepsilon_\text{MT}\) be this efficiency parameter.
We then say that this material is accreted instantly at the last simple binary evolution step.

Below, I show this assumption and compare it to the \(R_\text{RL}\) grid that I simulated last week.
These all have \(M_\text{TPAGB,i} = 2.2\; M_\odot \), \(q=0.6\) and \(\varepsilon =0.25\).
\begin{figure}[ht]
	\centering
	\incplot{w31-simple-vs-detailed}
\end{figure}

Clearly, the two match very well, and the small differences are probably due to the final assumption that there is \emph{no} wind during RLOF and the small
differences between the orbital evolution of the detailed \emph{MESA} models and the simple binary models.

It is also very easy to simulate large binaries that miss RLOF.

Below, I show the effects of the mass transfer efficiency on the final abundances.

\begin{figure}[ht]
	\centering
	\incplot{w31-change-eps}
\end{figure}

Interestingly, it seems like it is very easy to create barium stars even \emph{without any} RLOF
accretion.
The enrichment for the purple line comes only from the wind accretion.

If we do have \emph{RLOF} accretion, we can get higher abundances for smaller binaries.
These binaries will also have smaller final periods.
It also seems like we can create barium stars in very large binaries, just using wind accretion.

Lastly, we see a quite sharp drop between the models that undergo RLOF and those that do not.
I \emph{think} this is due to the envelope mass remaining during the last thermal pulse.
If the star has, say \(0.5\; M_\odot \) of envelope remaining, this will be accreted with RLOF
for models that \emph{just} undergo RLOF, while this will \emph{not} be accreted with the same
efficiency for models that just barely miss RLOF.

To test this, we can look at different masses, since this will shift the amount of
envelope mass remaining during the last thermal pulse.

Here again, we have \(q=0.6\) and \(\varepsilon =0.25\).
\begin{figure}[ht]
	\centering
	\incplot{w31-change-m}
\end{figure}

Here, it looks like the \(1.6\; M_\odot \) model has much more mass during the last thermal pulse
then for example, the \(2.8\; M_\odot \) star.

I confirm this below.
\begin{figure}[ht]
	\marginfignote{0}{Clearly, when plotting the radius as a fraction of the envelope (and also, but less clearly, just as
		the envelope mass), we can see that the last thermal pulses for the various models happen
		at different remaining envelope fractions.}
	\marginfignote{10}{The \(1.6\; M_\odot \) model has relatively a lot more envelope mass remaining
		when the last pulse happens then the \(2.8\; M_\odot \) model.
		I think it then makes sense why
		we see the sharp drop when the models miss RLOF and why the hight of this drop can depend quite
		strongly on the mass.
		It is really an effect of how much mass is remaining during the last TP.
	}
	\centering
	\incplot{w31-env-mass-r}
\end{figure}

We can also change the initial mass ratio \(q\).
This should have some interesting effects.
In particular, the \citet{saladino2018, saladino2019} wind prescription is dependent on \(q\), \(q\) shifts
the Roche lobe radius in an non-trivial way.
A smaller \(q\) will mean that \(R_\text{RL}/a\) is
larger, which means for the same initial \(R_\text{RL}\), the initial separation \(a_\text{i}\)
is larger.

Lastly, the RLOF mass transfer is depends on the initial mass ratio due to the fraction
\(\Delta M_\text{acc} / M_\text{acc}\).
In this simple case, \(\Delta M_\text{acc}\) is constant,
while \(M_\text{acc} = q \cdot M_\text{TPAGB}\).
So for smaller \(q\), we expect larger abundances
during RLOF.

\nextpage

\begin{figure}[ht]
	\centering
	\incplot{w31-change-q}
\end{figure}

We indeed see that the dependence on \(q\) is quite important.
In this simple case, a smaller
\(q\) will lead to \emph{higher} abundances during RLOF, but \emph{lower} abundances during the
wind mass loss regime.

In reality, the effect is probably more nuanced.
I would not at all be surprised
if there is some effect of \(q\) on the mass transfer efficiency.
Since models with lower
\(q\) tend to experience much more extreme RLOF, with higher mass loss rates, it would not
surprise me if more of this mass is lost, which would counteract the effect of higher abundances
due to the lower accretor mass.

We also see that the discontinueity between the RLOF and the wind regime is much more pronounced
for the lower \(q\) models.
I think this is probably because the mass transfer effiency of
wind is lowered with lower \(q\), while this is not the case during RLOF.

Let me explicitly test this.
Below, I show the amount of accreted mass for the different \(q\) models.
\begin{figure}[ht]
	\marginfignote{0}{We can indeed clearly see that the models with lower \(q\) experience
		much less wind accretion.}
	\marginfignote{8}{There also exists a nice powerlaw between the accreted material due
		to wind and the initial Roche lobe radius (or initial binary separation).}
	\centering
	\incplot{w31-change-q-macc}
\end{figure}

Since it is \emph{much} cheaper to compute the abundances of binary models using this simple approach
then using detailed \emph{MESA} simulations, it makes sense to explore the parameter space using
these simple models to get an \emph{estimate} for the abundances of the accretor.

I decided to use the following parameters:
\begin{itemize}
	\item  48 logspaced initial Roche lobe radii from \(125\; R_\odot \) to \(1250\; R_\odot \).
	\item  21 linearly spaced initial TPAGB masses from \(1.0\; M_\odot \) to \(3.0\; M_\odot \).
	\item 19 linearly spaced initial mass ratios from \(q = 0.1\) to \(1.0\).
	\item 20 logspaced RLOF accretion efficiencies from \(0.001\) to \(1\).
\end{itemize}

In total, this gives a grid of 383040 models, which only took a couple of hours to simulate\sidenote{Also because I used multithreading.
}
.

Obviously, because this is a 4D grid, it is impossible to show in its entirety at once.
To start, let me fix \(\varepsilon \) to \(\varepsilon \approx 0.113\).

\begin{figure}[ht]
	\marginfignote{4}{Fixed \(\varepsilon \approx 0.113\).}
	\marginfignote{6}{This grid shows the initial Roche lobe radius versus the initial TPAGB age in every panel.
		The different panels have different initial mass-ratios.
		The colors represent the expected final observed [s/Fe] for the accretor.

		We can see clear differences between the different panels, but (as expected), the minimum mass for significant s-process enrichment is independent of \(q\).
		It seems like we need at least a TPAGB star that is \(1.3\; M_\odot \) to get significantly large abundances.
		We then also see that the highest abundances occur in systems with the heaviest donors.

		This makes sense, these systems can donate more mass, \emph{and} they have more dredge-up then systems with lighter TPAGB donors.
		Lastly, we see that the initial mass ratio changes the \emph{width} of the enriched region quite significantly.
		This has two reasons:
		\begin{enumd}
			\item  The smaller \(q\) leads to faster orbital shrinkage due to winds, which shifts the boundary between RLOF and wind-only MT to lower initial Roche lobe radii.
			\item  Wind mass transfer is much less efficient for low \(q\) than it is for high \(q\).
			This means enrichment drops off much quicker in the wind-only regime for systems with small \(q\) then for systems with large \(q\).
		\end{enumd}
	}
	\centering
	\incplot{w31-R-M-s-q-panels}
\end{figure}

\begin{figure}[ht]
	\marginfignote{4}{Fixed \(\varepsilon \approx 0.113\).}
	\marginfignote{6}{This grid shows the initial Roche lobe radius versus the initial TPAGB age in every panel.
		The different panels have different initial mass-ratios.
		The colors represent the expected final observed [Ba/Fe] for the accretor.
	}
	\marginfignote{16}{There is only a small difference between these plots and the ones for [s/Fe].
		In general [Ba/Fe] is a bit larger, only because of the 7 most important \(s\)-process elements,
		barium is slightly more abundant than average.
	}
	\centering
	\incplot{w31-R-M-ba-q-panels}
\end{figure}

\begin{figure}[ht]
	\marginfignote{4}{Fixed \(\varepsilon \approx 0.113\).}
	\marginfignote{6}{This grid shows the initial Roche lobe radius versus the initial TPAGB age in every panel.
		The different panels have different initial mass-ratios.
		The colors represent the expected final observed [hs/ls] for the accretor.
	}
	\marginfignote{16}{I do not think this plot is very important.
		It seems to mostly be linked with the general [s/Fe] abundances.
		If I would have to say something about this, it \emph{seems} like the get slightly lower [hs/ls] for the most massive stars compared to the [s/Fe] abundances, which do not show this decrease with mass.
		This is consistent with \citet[][Figure 17]{karakas2016} where [hs/ls] rises to about 0.6 from \(M = 1.3\; M_\odot \) to \(\sim 2.4\; M_\odot \).
		Then it decreases to about \(0.3\) for \(M=3.0\; M_\odot \).
		Now since I am interpolating between these models, and slightly lower metallicity ones, we expect roughly the same properties, although maybe slightly higher because [hs/ls] seems to grow quite significantly with decreasing metallicity.
	}
	\centering
	\incplot{w31-R-M-hsls-q-panels}
\end{figure}

\nextpage
\begin{figure}[ht]
	\marginfignote{4}{Fixed \(\varepsilon \approx 0.113\).}
	\marginfignote{6}{This grid shows the initial Roche lobe radius versus the initial TPAGB age in every panel.
		The different panels have different initial mass-ratios.
		The colors represent the expected final \(\Delta M_\text{acc}\) for the accretor.
	}
	\marginfignote{16}{These panels very clearly show the boundary between the RLOF (+ wind) regime and the winds only regime.
		This is so visible because the winds, at least for the very large roche lobe radii, have a much lower accretion efficiency than the fixed value of \(0.113\) that I am plotting for RLOF.
		However, this is not the full picture.

		We can also see that for the larger \(q\)'s, there are 3 visible regimes.
		The lowest accretion regime is still the winds only regime, but the regime to the left of that now clearly shows
		the divide between systems that undergo significant amounts of wind before undergoing RLOF and systems that do not.
		The systems that \emph{do} experience a lot of winds will generally have
		slightly \emph{more} accreted material.
		At first, this might seem at odds with the lower accretion in the winds only regime, but it is actually not, because the accretion efficiency of the
		winds strongly depends on the size of the orbit.
		The much smaller orbits during this wind RLOF + wind regime can lead to \emph{larger} wind accretion efficiencies than the fixed \(0.113\) for RLOF.
	}
	\centering
	\incplot{w31-R-M-m2-q-panels}
\end{figure}

\begin{figure}[ht]
	\marginfignote{4}{Fixed \(\varepsilon \approx 0.113\).}
	\marginfignote{6}{This grid shows the initial Roche lobe radius versus the initial mass ratio \(q\) in every panel.
		The different panels have different initial TPAGB masses.
		The colors represent the expected final observed [s/Fe] values.
	}
	\marginfignote{16}{One thing that is immediately clear is that the lower mass models, up to say \(1.3\; M_\odot \) can \emph{not} or only under very specific circumstances create a barium star.
		It then also seems like once you reach a certain TPAGB mass, say something like \(2.2\; M_\odot \), the final s-process abundances of the accretor do not change a lot.
	}
	\marginfignote{30}{The increase in [s/Fe] towards lower \(q\) in the RLOF regime is due to the lower accretor mass, while the decrease in width is due to the less efficient wind mass transfer.
		This creates the general South-America shaped regions.
	}
	\marginfignote{40}{Lastly, there are quite clear horizontal lines visible in the abundances at
		the wind only regime.
		They seem to shift to lower \(q\) for larger donor masses.
		It took me a \emph{long} time to realise this is due to the ``full mixing'' that gets enabled for companion stars
		below \(0.8\; M_\odot \).
	}
	\centering
	\incplot{w31-R-q-s-M-panels}
\end{figure}

\begin{figure}[ht]
	\marginfignote{4}{Fixed mass \(M_\text{TPAGB,i} = 2.2\; M_\odot \).}
	\marginfignote{6}{This grid shows the initial Roche lobe radius versus the initial mass ratio \(q\) in every panel.
		The different panels have different \(\varepsilon_\text{MT}\) values during RLOF.
		The color represents the final [s/Fe] abundance.
	}
	\marginfignote{16}{We quite clearly see that, even for near-zero RLOF accretion efficiency, a large part of the parameter space can still create barium stars.
		This is quite surprising to me!
		Wind accretion alone (even in the RLOF regime) seems to be enough to create mild s-process enriched stars already.
	}
	\marginfignote{30}{Obviously, as we move to larger accretion efficiencies, the enhancements increase as well.
		There is however not \emph{much} difference between no RLOF accretion and something like \(\varepsilon_\text{RL} = 0.234\).
	}
	\marginfignote{40}{A last important difference is on the left sides of the panels.
		This regime consists of systems that only experience RLOF with \emph{very little to no} wind mass transfer.
		Therefore, a RLOF accretion efficiency of near-zero would just mean no mass is transferred to the accretor at all.
		This shifts the region of possible s-process enhanced stars from the \(M_\text{DUP}\) region to the RLOF + wind region.
	}
	\centering
	\incplot{w31-R-q-s-eps-panels}
\end{figure}

\begin{figure}[ht]
	\marginfignote{4}{Fixed mass \(M_\text{TPAGB,i} = 2.2\; M_\odot \).}
	\marginfignote{6}{This grid shows the initial Roche lobe radius versus the initial mass ratio \(q\) in every panel.
		The different panels have different \(\varepsilon_\text{MT}\) values during RLOF.
		The color represents the amount of accreted mass.
	}
	\marginfignote{16}{This plot makes the effect I was describing above visible.
		We can see that, for
		low accretion efficiencies, there is an island of systems that can get mass strictly due to wind,
		while for higher accretion efficiencies, the parameter space is divided into 2 where there is the
		RLOF (+ wind) regime and the wind-only regime.
	}
	\centering
	\incplot{w31-R-q-dm-eps-panels}
\end{figure}

\begin{figure}[ht]
	\marginfignote{4}{Fixed mass \(M_\text{TPAGB,i} = 1.3\; M_\odot \).}
	\marginfignote{6}{This grid shows the initial Roche lobe radius versus the initial mass ratio \(q\) in every panel.
		The different panels have different \(\varepsilon_\text{MT}\) values during RLOF.
		The color represents the final [s/Fe] abundance.
	}
	\marginfignote{16}{Since these panels are all for a much lower mass star, we can see that it is
		now much more difficult to create an s-process enhanced star.
		The value of \(\varepsilon \) is
		also not very dependent on the shape of the region again.
		We only really see major differences
		in a very small part of the parameter space, at Roche lobe radii that do not experience any wind, \emph{but} have already achieved some envelope enrichment.
	}
	\centering
	\incplot{w31-R-q-s-eps-2-panels}
\end{figure}

\bibliography{master.bib}
\end{document}
